% Copyright (c) 2026 Geovana Neves. Licensed under CC BY 4.0 (https://creativecommons.org/licenses/by/4.0/). See LICENSE-AND-AUTHORSHIP.md.
%
% 06-fsi/main.tex
%
% Guide 6 of the pyflightstream guides: Fluid-structure interaction. Built by ../build-all.ps1
% (or build-all.sh) with this folder and ../shared/ as include directories.
%
% No em dashes and no en dashes. No underscores in frame or section titles
% outside \code: beamer writes titles to .nav and .toc.

\documentclass[aspectratio=169,11pt]{beamer}

\input{preamble-slides.tex}
\input{hooks-slides.tex}
\ftsinput{boxes}
\ftsinput{hooks-contract}

\usepackage[nohyphen]{underscore}

\input{info.tex}

% The source line at the foot of a frame, as the gui deck's \guisource.
\newcommand{\fsisource}[1]{%
  \par\vfill
  {\tiny\color{SoftGray} Source: \pkgname{} [1]; #1.\par}}

% The source line of a frame that describes this version's behaviour from its
% requirements, written while the code that implements it was being merged.
\newcommand{\fsidecided}[1]{%
  \par\vfill
  {\tiny\color{SoftGray} Source: the requirements of this version; #1.\par}}

% Table columns, ragged right so a narrow cell never hyphenates or stretches:
% P{width} for a fixed column, R for the one that takes the rest.
\newcolumntype{P}[1]{>{\RaggedRight\arraybackslash}p{#1}}
\newcolumntype{R}{>{\RaggedRight\arraybackslash}X}

% A refusal as the package words it, quoted. Upright, because the italic
% underscore of the T1 sans face reads as a gap inside an identifier.
% \leavevmode first: a \color at the head of a p cell would open a blank line.
\newcommand{\says}[1]{{\leavevmode\color{DeepBlue!75!black}``#1''}}

% The clamp's hatching in the beam figure.
\usetikzlibrary{patterns}

% TikZ styles shared by the deck's figures.
\tikzset{
  fsibox/.style={draw=CourseBlue, fill=CourseBlue!6, rounded corners=1.5mm,
    align=center, font=\scriptsize, inner sep=1.6mm},
  fsifile/.style={draw=CourseGold, fill=CourseGold!12, align=center,
    font=\scriptsize\ttfamily, inner sep=1.2mm},
  fsiarrow/.style={-{Stealth[length=2mm]}, thick, CourseBlue},
  fsinote/.style={font=\scriptsize\itshape, text=SoftGray, align=center},
}

\begin{document}

\guidetitleframe{6}{Fluid-structure interaction}{FSI}
  {Fluid-structure interaction, as pyfs does it}
  {Sources: the package page ``FSI in a workspace'', the FSI tutorial and the FSI examples [1].}

% Copyright (c) 2026 Geovana Neves. Licensed under CC BY 4.0 (https://creativecommons.org/licenses/by/4.0/). See LICENSE-AND-AUTHORSHIP.md.
%
% fsi/text/01-reading.tex

\begin{frame}{What FSI means here, and how to read this deck}
\begin{columns}[T]
\begin{column}{0.55\textwidth}
\begin{itemize}\small
  \item \textbf{FSI} couples the aerodynamic loads on a blade or a wing to
        its structural deflection: FlightStream solves the flow, \pkgname{}
        solves the structure, and the solver deforms its surface with the
        displacements it is handed.
  \item The structure is \textbf{one Euler beam per blade}, or one along the
        wing (PyNite [5]), in flap and torsion, solved quasi-statically at
        every coupling call.
  \item A matrix row selects its structure by name, \code{FSI: f001}, and
        calibrates it from the matrix.
  \item Each frame names its source at its foot: the examples' numbers are
        the package's, on synthetic blades; the solver's behaviour is measured
        on 26.124 [2].
\end{itemize}
\end{column}
\begin{column}{0.41\textwidth}
\begin{block}{The optional extra}
\small The beam needs PyNite, from the \code{[fsi]} extra:
\code{pip install pyflightstream[fsi]}. Without it the package says how to
install it.
\end{block}
\begin{takeaway}
FSI runs on \code{steady}, \code{qsteady\_rotor} (sector) and \code{unsteady}
rows. A turning rotor, \code{unsteady\_rotor}, is refused: that coupling is
still being debugged.
\end{takeaway}
\end{column}
\end{columns}
\fsisource{FSI tutorial (\code{src/pyflightstream/fsi/README.md});
  \code{fsi/beam.py}; \code{extras.py}; the page ``FSI in a workspace''}
\end{frame}

% Re-checked against the merged FSI-1, FSI-GUARD and FSI-G code
% (cases/fsi_workspace.py FSI_WORKFLOW_STATE, validate_workspace_fsi,
% wire_fixed_wing_fsi). The qsteady_rotor sector row states the sector
% structural route that lands with the qsteady_rotor FSI wiring of this
% release.
\begin{frame}{Where FSI runs, and what the structure carries}
\relax % a body opening with a brace group would be read as the frame subtitle
{\ftstablesize\renewcommand{\arraystretch}{1.02}
\begin{tabularx}{\linewidth}{@{}P{0.20\linewidth}P{0.10\linewidth}R@{}}
\toprule
\textbf{Workflow} & \textbf{FSI} & \textbf{What the structure carries, and why} \\
\midrule
\code{steady} & allowed & a fixed wing, one clamped beam: the aerodynamic
  loads and the structure's \textbf{own weight}; nothing turns, so no
  centrifugal load. At most 50 coupling iterations \\
\code{qsteady\_rotor}, sector & allowed & one blade held still in a free stream
  turning at the rotor's speed: the structural solver applies the
  \textbf{centrifugal} loads and stiffening at the \textbf{free-stream
  rotation's} rpm, \(\Omega = |\mathrm{rpm}|\,\pi/30\) \\
\code{qsteady\_rotor}, wheel & refused & several steady clockings averaged; a
  blade deformed once per clocking is not the state of one structure \\
\code{unsteady} & allowed & a fixed wing, nothing turning: the loads and its
  own weight, one coupling call per time step \\
\code{unsteady\_rotor} & refused & \says{FSI on unsteady\_rotor is still in
  debug on this release}: on 26.124 the morph of a turning blade is applied
  to the blade at its imported azimuth [2] \\
\bottomrule
\end{tabularx}}
\par\vspace{1mm}
{\scriptsize A fixed wing's FSI input states \code{[config.wing]} and couples
on \code{steady} or \code{unsteady} only; a blade's input, without it, couples
on the rotor route only. A quasi-steady answer is valid only for an isolated,
axisymmetric rotor, and only while the flow at each station changes slowly
beside the revolution (Guide 1, the workflows; Guide 5).\par}
\fsisource{\code{cases/fsi\_workspace.py}, \code{FSI\_WORKFLOW\_STATE},
  \code{validate\_workspace\_fsi}, \code{STEADY\_AEROELASTIC\_ITERATIONS};
  \code{cases/workflows.py}; RPT-093 [2]}
\end{frame}


\begin{frame}{Contents}
\small
\begin{multicols}{2}
\tableofcontents
\end{multicols}
\end{frame}

% Copyright (c) 2026 Geovana Neves. Licensed under CC BY 4.0 (https://creativecommons.org/licenses/by/4.0/). See LICENSE-AND-AUTHORSHIP.md.
%
% fsi/text/02-the-loop.tex

\section{The coupling loop}

\begin{frame}{The loop, one coupling call}
\begin{center}
\begin{adjustbox}{max width=\linewidth, max totalheight=0.54\textheight}
\begin{tikzpicture}[node distance=5mm]
  \node[fsibox, minimum width=3.3cm, minimum height=1.35cm] (fs)
    {\textbf{FlightStream}\\a steady iteration or\\a time step, on the\\current mesh};
  \node[fsibox, right=9mm of fs, minimum width=4.4cm, minimum height=1.35cm] (post)
    {\textbf{post-processing script} \texttt{fsi\_post.txt}\\[0.5mm]
     \texttt{UPDATE\_ALL\_SURFACE\_SECTIONS}\\
     \texttt{COMPUTE\_SURFACE\_SECTIONAL\_LOADS NEWTONS}\\
     \texttt{EXPORT\_SURFACE\_SECTIONAL\_LOADS}};
  \node[fsifile, right=9mm of post, minimum height=1.35cm, text width=3.7cm] (loads)
    {FS\_SurfaceSection\_Loads.txt\\[0.5mm]
     {\rmfamily\itshape per section: Fx, Fz, moment about the quarter chord}};
  \node[fsibox, below=15mm of loads, minimum width=3.9cm, minimum height=1.7cm,
        text width=3.8cm] (step)
    {\textbf{\texttt{pyfs-fsi} step}\\[0.5mm]
     loads moved to the elastic axis\\
     one beam solve per blade\\
     relaxed against the last call\\
     twist written as node translations};
  \node[fsifile, left=12mm of step, minimum height=1.35cm, text width=2.7cm] (disp)
    {FSIDisp.txt\\[0.5mm]
     {\rmfamily\itshape one translation per node, in import order}};
  \node[fsibox, below=15mm of fs, minimum width=3.3cm, minimum height=1.35cm] (morph)
    {the solver moves\\the blade nodes\\and deforms the mesh};
  \node[fsifile, right=6mm of step, yshift=4mm] (cfg) {config.json};
  \node[fsifile, below=2mm of cfg] (state) {state.json};

  \draw[fsiarrow] (fs) -- (post);
  \draw[fsiarrow] (post) -- (loads);
  \draw[fsiarrow] (loads) -- (step);
  \draw[fsiarrow] (step) -- (disp);
  \draw[fsiarrow] (disp) -- (morph);
  \draw[fsiarrow] (morph) -- node[fsinote, left] {next\\call} (fs);
  \draw[fsiarrow, CourseGold] (cfg.west) -- (cfg.west -| step.east);
  \draw[fsiarrow, CourseGold, {Stealth[length=2mm]}-{Stealth[length=2mm]}]
    (state.west) -- (state.west -| step.east);
\end{tikzpicture}
\end{adjustbox}
\end{center}
\vspace{-1mm}
{\small FlightStream's Aeroelastic Toolbox calls the executable once per
aeroelastic iteration of a steady row, or once per time step of an unsteady
one (\code{SET\_AEROELASTIC\_ITERATIONS 1}). Each call is stateless: what
persists is in \code{state.json}. Loads and displacements are in each blade's
own frame, or the wing's; the package never sees azimuth.\par}
\fsisource{FSI tutorial, ``The loop in one paragraph'', ``cli'';
  \code{cases/fsi\_workspace.py}, \code{wire\_workspace\_fsi};
  \code{fsi/\_\_init\_\_.py}}
\end{frame}

\begin{frame}{The structure: one Euler beam per blade, or along the wing}
\begin{columns}[T]
\begin{column}{0.52\textwidth}
\begin{itemize}\small
  \item One node per radial station, on the elastic axis; one member per
        bay, with the bay-averaged EI and GJ; the root station clamped.
  \item Unit elastic moduli, \(E = G = 1\): the section constants handed to
        PyNite are the configuration's EI and GJ, so no modulus is applied
        twice.
  \item Flap deflection \(w\) and elastic twist \(\theta\) carry mass
        (\(\mu\,l\) and \((I_1 + I_2)\,l\)); every other degree of freedom is
        condensed out of the eigenproblem exactly.
  \item Static solves are linear, or P-Delta whenever axial tension is
        present [6, 7].
\end{itemize}
\end{column}
\begin{column}{0.44\textwidth}
\begin{center}
\begin{adjustbox}{max width=\linewidth}
\begin{tikzpicture}[x=1cm,y=1cm]
  % the clamp
  \fill[pattern=north east lines, pattern color=SoftGray] (-0.35,-0.8) rectangle (0,0.8);
  \draw[thick, SoftGray] (0,-0.8) -- (0,0.8);
  % the beam and its stations
  \draw[very thick, CourseBlue] (0,0) -- (5.6,0);
  \foreach \x/\n in {0/0,1.4/1,2.8/2,4.2/3,5.6/4} {
    \fill[CourseBlue] (\x,0) circle (1.6pt);
  }
  \node[font=\tiny, below, text=SoftGray] at (0.7,0) {bay};
  \draw[SoftGray, decorate, decoration={brace, amplitude=3pt, mirror}]
    (1.4,-0.28) -- (2.8,-0.28) node[midway, below=2pt, font=\tiny] {EI, GJ averaged};
  % flap and twist at the tip
  \draw[-{Stealth[length=1.6mm]}, thick, CourseRed] (5.6,0) -- (5.6,0.9)
    node[above, font=\scriptsize] {\(w\), flap};
  \draw[-{Stealth[length=1.6mm]}, thick, CourseRed]
    (6.05,-0.35) arc[start angle=-60, end angle=200, x radius=0.25, y radius=0.45];
  \node[font=\scriptsize, CourseRed, right] at (6.2,0.1) {\(\theta\), twist};
  % the tension
  \draw[-{Stealth[length=1.6mm]}, thick, CourseGold] (3.0,0.45) -- (4.2,0.45)
    node[midway, above, font=\scriptsize, text=black] {\(\mu\,\Omega^2 r\)};
  \node[font=\tiny, text=SoftGray, align=center] at (-0.2,-1.15) {root\\clamped};
  \node[font=\tiny, text=SoftGray] at (5.6,-0.95) {tip};
\end{tikzpicture}
\end{adjustbox}
\end{center}
\begin{takeaway}
The stiffnesses in \code{config.json} are the beam's stiffnesses, typed or
generated, exactly as written.
\end{takeaway}
\end{column}
\end{columns}
\fsisource{FSI tutorial, ``beam''; \code{fsi/beam.py} module docstring}
\end{frame}

\begin{frame}{What rotation adds}
\begin{itemize}\small
  \item \textbf{Tension.} The distributed axial force \(\mu\,\Omega^2 r\)
        makes the tension \(N(r)\) emerge in the solver, and P-Delta turns it
        into bending stiffness, with no manual correction terms: the
        centrifugal stiffening of a rotating blade [8, 9, 10].
  \item \textbf{Propeller moment.} \(-\Omega^2 (I_1 - I_2)\sin\theta\cos\theta\)
        drives each section toward flat pitch. It depends on the total pitch,
        geometric plus elastic, so each call re-solves the beam a few times
        until the twist stops moving [9, 12].
  \item \textbf{In-plane softening.} The part of the flap deflection that
        lies in the plane of rotation feels the centrifugal field pulling it
        outward, a term \(\mu\,\Omega^2 \sin^2\beta\,w\), \(\beta\) the
        section's pitch, that lowers the stiffness; it is applied with the
        tension [8, 9].
  \item \textbf{Not modelled: the trapeze effect.} A declared limitation;
        omitting it overpredicts the elastic twist, which is conservative.
\end{itemize}
\fsisource{FSI tutorial, ``The loop in one paragraph'' and ``centrifugal'';
  \code{fsi/centrifugal.py} module docstring}
\end{frame}

\begin{frame}{A fixed wing: its own weight, and where the model stops}
\begin{itemize}\small
  \item \textbf{A wing that does not turn} has no centrifugal load. The
        structural solver applies the structure's \textbf{own weight}, gravity
        on its mass distribution, beside the aerodynamic loads, on a
        \code{steady} or an \code{unsteady} row; the loop is the same.
  \item \textbf{A blade held still} in a turning free stream
        (\code{qsteady\_rotor}, sector) is the opposite case: the aerodynamics
        see no motion, and the structural solver still applies the
        centrifugal loads at the row's rotation speed.
  \item \textbf{Quasi-static.} The mean deformation is exact within the beam
        idealisation; the once-per-revolution content is trustworthy only well
        below resonance, roughly \(n\,\Omega / \omega_n \le 0.3\) [12, 13].
        The convergence log prints this boundary in its header.
\end{itemize}
\begin{alertblock}{The formulas are transcribed, not re-verified}
\small The primary sources of the model have not been independently checked
against the implemented formulas; each formula lives in one small function
that cites its source, so a correction stays local.
\end{alertblock}
\fsidecided{\code{fsi/centrifugal.py}; the page ``FSI in a workspace''}
\end{frame}

\begin{frame}{Twist travels as three nodes per station}
\begin{center}
\begin{adjustbox}{max width=0.62\linewidth, max totalheight=0.36\textheight}
\begin{tikzpicture}[x=1cm,y=1cm]
  % undeformed section
  \draw[SoftGray, dashed] (0,0) ellipse (2.4 and 0.32);
  \draw[SoftGray] (-2.9,0) -- (2.9,0);
  \foreach \x in {-1.2,0,1.2} \fill[SoftGray] (\x,0) circle (2pt);
  % deformed section: translated by w and turned by theta
  \begin{scope}[shift={(0,1.15)}, rotate=9]
    \draw[CourseBlue, thick] (0,0) ellipse (2.4 and 0.32);
    \fill[CourseBlue] (-1.2,0) circle (2.2pt);
    \fill[CourseRed] (0,0) circle (2.2pt);
    \fill[CourseBlue] (1.2,0) circle (2.2pt);
  \end{scope}
  \draw[-{Stealth[length=1.6mm]}, CourseRed, thick] (0,0) -- (0,1.12)
    node[midway, left, font=\scriptsize] {\(w\)};
  \draw[-{Stealth[length=1.4mm]}, CourseBlue] (1.2,0) -- (1.2,1.33)
    node[midway, right, font=\scriptsize] {\(w + \theta d\)};
  \draw[-{Stealth[length=1.4mm]}, CourseBlue] (-1.2,0) -- (-1.2,0.97)
    node[midway, left, font=\scriptsize] {\(w - \theta d\)};
  \node[font=\tiny, text=SoftGray] at (1.2,-0.55) {toward LE};
  \node[font=\tiny, text=SoftGray] at (0,-0.55) {elastic axis};
  \node[font=\tiny, text=SoftGray] at (-1.2,-0.55) {toward TE};
  \draw[SoftGray, |-|] (0,-0.9) -- (1.2,-0.9) node[midway, below, font=\tiny] {\(d = 0.25\,c\)};
\end{tikzpicture}
\end{adjustbox}
\end{center}
\vspace{-1mm}
\begin{columns}[T]
\begin{column}{0.48\textwidth}
\begin{itemize}\footnotesize
  \item \code{FSIDisp.txt} carries translations only, so twist is encoded by
        three nodes per station, \(d = 0.25\,c\) apart
        (\code{node\_offset\_chord\_fraction}).
  \item Rigid section kinematics, \(\delta y = w + \theta\,d\): linear, so the
        inverse is exact to machine precision.
\end{itemize}
\end{column}
\begin{column}{0.48\textwidth}
\begin{itemize}\footnotesize
  \item One generator makes the node file, the ordering map and the
        \code{FSIDisp.txt} rows: the solver applies them strictly in import
        order, and a drift would corrupt the geometry in silence.
  \item One node file serves every blade, imported once per blade frame, in
        blade order.
\end{itemize}
\end{column}
\end{columns}
\fsisource{FSI tutorial, ``kinematics'' and ``nodes'';
  \code{fsi/nodes.py} module docstring; \code{fsi/config.py},
  \code{FsiConfig.node\_offset\_chord\_fraction}}
\end{frame}

\begin{frame}{The driver: four phases, keyed on the step counter}
\relax % a body opening with a brace group would be read as the frame subtitle
{\ftstablesize\renewcommand{\arraystretch}{1.2}
\begin{tabularx}{\linewidth}{@{}P{0.23\linewidth}R@{}}
\toprule
\textbf{Phase} & \textbf{Behaviour, and its key in \code{config.json} \code{phases}} \\
\midrule
1, wake development & zero displacements on the rigid blade;
  \code{wake\_development\_revolutions} (1.0) \\
2, averaged coupling & loads averaged over the window, relaxed updates
  \(d_{new} = d_{old} + \lambda (d_{calc} - d_{old})\);
  \code{averaging\_window\_revolutions} (1.0), \code{coupling\_relaxation} (0.4) \\
3, convergence watch & as 2, and converged when BOTH hold across consecutive
  revolutions: tip twist change below \code{tip\_twist\_tolerance\_deg} (0.05)
  and relative change of the integrated normal force below
  \code{thrust\_tolerance\_fraction} (0.02) \\
4, recording & instantaneous loads, \(\lambda = 1\) by design, twist recorded
  per step; \code{recording\_revolutions} (1.0) \\
\bottomrule
\end{tabularx}}
\par\vspace{1.5mm}
\begin{itemize}\scriptsize
  \item \textbf{Freshness is asserted}: each call must see an advancing solver
        iteration in the loads header; a repeat is refused, since it means more
        than one FSI iteration per time step.
  \item \textbf{Frozen mode}: a \code{fsi\_frozen\_displacements.txt} in the
        run folder is replayed verbatim at every call, with no loads and no solve.
\end{itemize}
\fsisource{FSI tutorial, ``driver and state''; \code{fsi/config.py},
  \code{PhaseSchedule} (defaults in brackets)}
\end{frame}

\begin{frame}{The structural nodes, and how the surface follows them}
\begin{columns}[T]
\begin{column}{0.52\textwidth}
\begin{itemize}\small
  \item \textbf{Inside the component.} The nodes lie on the camber line of
        each station, turned by its local twist, inside the blade or the
        wing; a node outside is refused at plan.
  \item \textbf{In the frame that turns with the blade}, so a displacement
        means the same thing at every position.
  \item \textbf{The surface by its boundary}: the aeroelastic surface is
        assigned by the boundary's own identifier in the mesh.
  \item \textbf{After the structural call}: a coupled point exports once the
        last call has moved the surface; a steady run waits for it.
\end{itemize}
\end{column}
\begin{column}{0.44\textwidth}
\begin{block}{A radial basis function carries the nodes to the panels}
\small The displacement known at the nodes is interpolated onto every
aerodynamic panel [14, 15, 16]. A beam is a \textbf{line} of nodes: a
compactly supported kernel reaches only the panels within its radius of that
line, and the rest of the surface barely moves. The coupled row uses the
\textbf{multiquadric} [14], whose support is global.
\end{block}
\end{column}
\end{columns}
\fsidecided{\code{fsi/nodes.py}; \code{cases/fsi\_workspace.py}; Guide 4,
  \code{aeroelastic\_rbf\_type}}
\end{frame}

\begin{frame}{Route B: the direct transform, a sanity check}
\begin{columns}[T]
\begin{column}{0.52\textwidth}
\begin{itemize}\small
  \item The same beam, solved on the solver's sectional loads, gives flap
        \(w(r)\) and twist \(\theta(r)\); \textbf{every mesh vertex is moved}
        about the elastic axis by that \(w\) and \(\theta\), and the deformed
        geometry is solved again, rigid. No interpolation, no in-run
        coupling.
  \item Repeated until the tip flap and twist stop moving, it is a
        partitioned static iteration: quasi-steady by construction.
  \item A zero deformation through the same chain reproduces the rigid
        answer, which is its own control.
\end{itemize}
\end{column}
\begin{column}{0.44\textwidth}
\begin{ftslimit}[What route B is for]
\small A simplified run, a check of the coupled route against an
independent one, and an FSI option where the coupled route is not
available. It is \textbf{not} the coupled route: it carries no dynamics and
no in-run deformation. The worked example of Part 7 was solved this way.
\end{ftslimit}
\end{column}
\end{columns}
\fsidecided{the worked example's route diagnostic [3]}
\end{frame}

% Copyright (c) 2026 Geovana Neves. Licensed under CC BY 4.0 (https://creativecommons.org/licenses/by/4.0/). See LICENSE-AND-AUTHORSHIP.md.
%
% fsi/text/03-the-input.tex

\section{The structural input}

\begin{frame}[fragile]{One structural input, selected by name}
\begin{columns}[T]
\begin{column}{0.50\textwidth}
\begin{itemize}\small
  \item One configuration per file, \code{inputs/fsi/f001.toml}: \code{f},
        then letters, digits, \code{\_} or \code{-}; no extension in the row.
  \item The row selects it in \code{VAR\_NAMES\_VALUES}, with any
        calibration beside it.
  \item The resolver attaches the \emph{effective} \code{FsiConfig} and its
        provenance to the case. No solver runs while loading or planning.
  \item The run writes both files in the point's working directory and
        records their hashes. The source is hashed too: changing it means
        planning again.
\end{itemize}
\end{column}
\begin{column}{0.46\textwidth}
\begin{lstlisting}[style=ftsbase, basicstyle=\ttfamily\tiny, breaklines=true]
FSI: f001 / FSI_BENDING_STIFFNESS_N_M2_FACTOR: 1.15
\end{lstlisting}
\vspace{1mm}
\begin{center}
\begin{adjustbox}{max width=\linewidth, max totalheight=0.36\textheight}
\begin{tikzpicture}
  \node[fsifile, minimum width=2.9cm] (toml) at (-1.8,0) {inputs/fsi/f001.toml};
  \node[fsifile, minimum width=2.9cm] (row) at (1.8,0) {the row's FSI keys};
  \node[fsibox, minimum width=4.6cm] (res) at (0,-1.15) {the resolver, at plan:\\base, factors, effective};
  \node[fsifile, minimum width=2.9cm] (cfg) at (-1.8,-2.3) {config.json};
  \node[fsifile, minimum width=2.9cm] (prov) at (1.8,-2.3) {fsi-provenance.json};
  \draw[fsiarrow] (toml.south) -- (res.north -| toml.south);
  \draw[fsiarrow] (row.south) -- (res.north -| row.south);
  \draw[fsiarrow] (res.south -| cfg.north) -- (cfg.north);
  \draw[fsiarrow] (res.south -| prov.north) -- (prov.north);
  \node[fsinote] at (0,-2.85) {written and hashed in the point's working directory};
\end{tikzpicture}
\end{adjustbox}
\end{center}
\begin{takeaway}
The geometry library is never a destination for these generated files.
\end{takeaway}
\end{column}
\end{columns}
\fsisource{\code{docs/fsi-workspace.md}, opening section;
  \code{workspace/fsi\_setup.py}, \code{resolve\_row\_fsi};
  \code{docs/srs/data-model.md}, the \code{inputs/} tree}
\end{frame}

\begin{frame}{Two ways to state the blade}
\relax % a body opening with a brace group would be read as the frame subtitle
{\ftstablesize\renewcommand{\arraystretch}{1.25}
\begin{tabularx}{\linewidth}{@{}P{0.20\linewidth}RR@{}}
\toprule
 & \textbf{\code{mode = "supplied"}} & \textbf{\code{mode = "calculated"}} \\
\midrule
What you write & the \code{FsiConfig} scalars under \code{[config]} and the
  complete \code{BladeProperties} arrays under \code{[config.blade]} &
  the \code{[config]} scalars, a \code{material}, and \code{[sections]}: one
  closed contour per station, in metres \\
What the package computes & nothing: the distributions are the file's, then
  calibrated & area, centroid, principal inertias, the torsion constant, then
  EI, GJ, running mass and mass inertias from the material \\
Where the numbers come from & you: the example distributions are synthetic,
  not a material or airframe dataset & the geometry and ONE sourced
  material, recorded with the numbers \\
Refused & a \code{material} or \code{[sections]} table; any material factor &
  a \code{[config.blade]} table; no material; no \code{[sections]} \\
\bottomrule
\end{tabularx}}
\par\vspace{2mm}
\begin{takeaway}
Either way the result is the existing solid homogeneous Euler beam: this input
adds no shell, spar, hollow-section, shear-centre or second structural model.
\end{takeaway}
\fsisource{\code{docs/fsi-workspace.md}, ``Supplied distributions'' and
  ``Calculated solid sections''; \code{workspace/fsi\_setup.py},
  \code{resolve\_fsi\_setup}}
\end{frame}

\begin{frame}[fragile]{Supplied distributions}
\begin{columns}[T]
\begin{column}{0.47\textwidth}
\begin{lstlisting}[style=ftsbase, basicstyle=\ttfamily\tiny]
mode = "supplied"

[calibration]
bending_stiffness_n_m2 = 1.05

[config]
blade_count = 2
omega_rad_per_s = 0.0

[config.blade]
station_radii_m = [0.1, 0.5]
chord_m = [0.04, 0.04]
mass_per_length_kg_per_m = [1.0, 2.0]
inertia_major_kg_m = [0.001, 0.002]
inertia_minor_kg_m = [0.0001, 0.0002]
bending_stiffness_n_m2 = [10.0, 20.0]
torsion_stiffness_n_m2 = [5.0, 10.0]
elastic_axis_offset_chordwise_m = [0.0, 0.0]
elastic_axis_offset_normal_m = [0.0, 0.0]
cg_offset_chordwise_m = [0.0, 0.0]
cg_offset_normal_m = [0.0, 0.0]
geometric_pitch_deg = [0.0, 0.0]
\end{lstlisting}
\end{column}
\begin{column}{0.49\textwidth}
{\ftstablesize\renewcommand{\arraystretch}{1.15}
\begin{tabular}{@{}ll@{}}
\toprule
\textbf{Array} & \textbf{Unit} \\
\midrule
\code{station\_radii\_m}, strictly increasing & m \\
\code{mass\_per\_length\_kg\_per\_m} & kg/m \\
\code{bending\_stiffness\_n\_m2} (EI, flap) & N\,m\(^2\) \\
\code{torsion\_stiffness\_n\_m2} (GJ) & N\,m\(^2\) \\
\code{inertia\_major\_kg\_m}, \code{\_minor\_kg\_m} & kg\,m \\
offsets, \code{chord\_m} & m \\
\bottomrule
\end{tabular}}
\par\vspace{1.5mm}
\begin{itemize}\scriptsize
  \item Every array has the length of \code{station\_radii\_m}; values are
        finite; mass, EI, GJ and chord positive; principal inertias not
        negative.
  \item With the row factor 1.15 of the previous part, the effective EI is
        \code{[11.5, 23.0]}: the row replaces the file's 1.05, it does not
        multiply it.
\end{itemize}
{\scriptsize\itshape A zero-speed input like this one is a valid structural
input for offline checks, not a runnable coupled rotor row.\par}
\end{column}
\end{columns}
\fsisource{\code{docs/fsi-workspace.md}, ``Supplied distributions'';
  \code{fsi/config.py}, \code{BladeProperties}}
\end{frame}

\begin{frame}[fragile]{Calculated solid sections}
\begin{columns}[T]
\begin{column}{0.50\textwidth}
\begin{lstlisting}[style=ftsbase, basicstyle=\ttfamily\tiny]
mode = "calculated"
material = "ti-6al-4v-grade5-annealed"

[config]
blade_count = 2
omega_rad_per_s = 0.0

[sections]
station_radii_m = [0.1, 0.5]
chord_m = [0.04, 0.04]
geometric_pitch_deg = [0.0, 0.0]
geometry_source = "Synthetic rectangular sections in SI units"
torsion_grid_cells = 16
sections_m = [
  [[-0.02, -0.002], [0.02, -0.002],
   [0.02, 0.002], [-0.02, 0.002]],
  [[-0.02, -0.002], [0.02, -0.002],
   [0.02, 0.002], [-0.02, 0.002]]
]
\end{lstlisting}
\end{column}
\begin{column}{0.46\textwidth}
\begin{itemize}\small
  \item \code{sections\_m}: one closed polygon per station, chordwise and
        normal metres about the pitch axis.
  \item \code{geometry\_source}: what the contours were built from, in
        words; it travels with the numbers.
  \item \code{torsion\_grid\_cells}: cells across the thickness for the
        torsion solve, 8 to 512, 64 when omitted.
  \item No \code{[config.blade]}: this mode derives it.
\end{itemize}
\begin{takeaway}
A 40 by 4 mm titanium strip; example 1 turns these contours into EI.
\end{takeaway}
\end{column}
\end{columns}
\fsisource{\code{docs/fsi-workspace.md}, ``Calculated solid sections'';
  \code{workspace/fsi\_setup.py}, \code{SolidSections}, \code{FSI\_TEMPLATE}}
\end{frame}

\begin{frame}{The material database: no number without a source}
\relax % a body opening with a brace group would be read as the frame subtitle
{\ftstablesize\renewcommand{\arraystretch}{1.2}
\begin{tabularx}{\linewidth}{@{}lrrrrR@{}}
\toprule
\textbf{Key} & \(\rho\) [kg/m\(^3\)] & \(E\) [GPa] & \(G\) [GPa] & \(\nu\) & \textbf{Source, as recorded} \\
\midrule
\code{ti-6al-4v-grade5-annealed} & 4430 & 113.8 & 44.0 & 0.342 &
  ASM Aerospace Specification Metals data sheet, annealed; \(G\) tabulated \\
\code{al-7075-t6} & 2810 & 71.7 & 26.9 & 0.33 &
  ASM Aerospace Specification Metals data sheet, T6 typical values; \(G\) tabulated \\
\bottomrule
\end{tabularx}}
\par\vspace{2mm}
\begin{itemize}\small
  \item Each entry's four numbers come from ONE cited data set, with its table,
        condition, address and reading date. Database version \code{1}; a
        changed number is a new version.
  \item Your own material: a \code{[material]} table stating every field, with
        a \code{[material.source]} of \code{document}, \code{table},
        \code{condition}, \code{url}, \code{consulted}. Density, \(E\), \(G\)
        finite and positive; \(-1 < \nu < 0.5\).
  \item Room-temperature typical values for a structural model of a blade, not
        the design allowables of a certified part.
\end{itemize}
\fsisource{\code{fsi/materials.py}; \code{docs/fsi-workspace.md},
  ``Calculated solid sections''; FSI tutorial, ``sections and materials''}
\end{frame}

\begin{frame}{The solid-section calculator}
\begin{columns}[T]
\begin{column}{0.60\textwidth}
{\ftstablesize\renewcommand{\arraystretch}{1.12}
\begin{tabularx}{\linewidth}{@{}P{0.50\linewidth}R@{}}
\toprule
\textbf{Field} & \textbf{Generated as} \\
\midrule
\code{mass\_per\_length\_kg\_per\_m} & \(\rho A\) \\
\code{inertia\_major\_kg\_m}, \code{\_minor} & \(\rho\) times the larger, smaller principal second moment \\
\code{bending\_stiffness\_n\_m2} & \(E\) times the second moment about the chordwise centroidal axis \\
\code{torsion\_stiffness\_n\_m2} & \(G J\), \(J\) computed numerically \\
\code{elastic\_axis\_offset\_*} & the centroid's position from the pitch axis \\
\code{cg\_offset\_*} & zero \\
\bottomrule
\end{tabularx}}
\end{column}
\begin{column}{0.36\textwidth}
\begin{itemize}\scriptsize
  \item Area, centroid, second moments: exact closed forms of the polygon.
        A contour that crosses or touches itself, repeats a vertex or folds
        back is refused by name.
  \item \(J\): the Prandtl stress function by finite differences cut exactly at
        the polygon [6, 7]. 16 cells across are within about 1\,\% of the closed
        forms, 32 within 0.3\,\%, 64 within 0.08\,\%.
  \item The thin-section \(\frac13\!\int t^3\,ds\) is kept as a cross-check,
        never as the value.
\end{itemize}
\end{column}
\end{columns}
\par\vspace{1mm}
\begin{alertblock}{Stated hypotheses, recorded with every generated blade}
\scriptsize SOLID, HOMOGENEOUS sections; the elastic axis at the CENTROID, the
shear centre not computed; a zero elastic-axis-to-CG offset, so no bend-twist
coupling from it; flap bending about the chordwise axis.
\end{alertblock}
\fsisource{FSI tutorial, ``sections and materials''}
\end{frame}

% Copyright (c) 2026 Geovana Neves. Licensed under CC BY 4.0 (https://creativecommons.org/licenses/by/4.0/). See LICENSE-AND-AUTHORSHIP.md.
%
% fsi/text/04-calibration.tex

\section{Calibration and provenance}

\begin{frame}{Calibration factors: unity, the file, the matrix}
\begin{columns}[T]
\begin{column}{0.60\textwidth}
{\ftstablesize\renewcommand{\arraystretch}{1.2}
\begin{tabularx}{\linewidth}{@{}P{0.27\linewidth}R@{}}
\toprule
\textbf{Level} & \textbf{Property names} \\
\midrule
Distributions, either mode & \code{mass\_per\_length\_kg\_per\_m},
  \code{inertia\_major\_kg\_m}, \code{inertia\_minor\_kg\_m},
  \code{bending\_stiffness\_n\_m2}, \code{torsion\_stiffness\_n\_m2} \\
Orientation & \code{geometric\_pitch\_deg} \\
Offsets & \code{elastic\_axis\_offset\_chordwise\_m}, \code{\_normal\_m},
  \code{cg\_offset\_chordwise\_m}, \code{\_normal\_m} \\
Material, calculated mode only & \code{density\_kg\_per\_m3},
  \code{youngs\_modulus\_pa}, \code{shear\_modulus\_pa}, \code{poisson\_ratio} \\
\bottomrule
\end{tabularx}}
\par\vspace{1.5mm}
{\small The matrix key is \code{FSI\_} + the name in capitals + \code{\_FACTOR}:
\code{FSI\_TORSION\_STIFFNESS\_N\_M2\_FACTOR}.\par}
\end{column}
\begin{column}{0.36\textwidth}
\begin{center}
\begin{adjustbox}{max width=\linewidth}
\begin{tikzpicture}[node distance=3mm]
  \node[fsibox, minimum width=4.2cm] (m) {\textbf{matrix}: \code{FSI\_*\_FACTOR}\\on the row, even \code{1.0}};
  \node[fsibox, below=of m, minimum width=4.2cm, fill=CourseGold!10, draw=CourseGold] (f) {\textbf{file}: the TOML's\\\code{[calibration]} table};
  \node[fsibox, below=of f, minimum width=4.2cm, fill=SoftGray!10, draw=SoftGray] (u) {\textbf{unity}: stated nowhere};
  \draw[fsiarrow] ([xshift=3mm]u.south east) -- ([xshift=3mm]m.north east)
    node[midway, fsinote, rotate=90, above] {takes precedence};
\end{tikzpicture}
\end{adjustbox}
\end{center}
\begin{itemize}\scriptsize
  \item Every factor is dimensionless, finite and positive.
  \item An explicit matrix value replaces the file factor exactly once; it
        is never multiplied by it.
\end{itemize}
\end{column}
\end{columns}
\fsisource{\code{docs/fsi-workspace.md}, ``Calibration and provenance'';
  \code{\_fsi\_calibration.py}, \code{MATRIX\_FACTORS}}
\end{frame}

\begin{frame}{One property, scaled once}
\begin{columns}[T]
\begin{column}{0.54\textwidth}
{\small Calibrating a source and a property derived from it would scale the
same number twice. Refused, in calculated mode:\par}
{\ftstablesize\renewcommand{\arraystretch}{1.2}
\begin{tabularx}{\linewidth}{@{}P{0.46\linewidth}R@{}}
\toprule
\textbf{Source factor} & \textbf{with any of} \\
\midrule
density & mass per length, either mass inertia \\
Young's modulus \(E\) & EI \\
shear modulus \(G\) & GJ \\
\(E\) or \(\nu\), when \(G\) is derived & a separate \(G\) or GJ factor \\
\bottomrule
\end{tabularx}}
\par\vspace{1.5mm}
{\scriptsize \says{\ldots: calibrating youngs\_modulus\_pa and its derived
bending\_stiffness\_n\_m2 would scale the same property twice; choose one
calibration level.}\par}
\end{column}
\begin{column}{0.42\textwidth}
\begin{itemize}\scriptsize
  \item The legacy \code{config.stiffness\_scale\_factor} must be 1.0 in a
        workspace input: calibrate EI and GJ explicitly instead of a second
        scaling at solve time.
  \item Supplied mode takes no material factor: calibrate the distribution
        itself.
  \item Every effective property is validated again, including major
        principal inertia not below minor.
\end{itemize}
\begin{ftsevidence}[File 1.05, row 1.15]
\scriptsize Effective EI = base \(\times\) 1.15, not base \(\times\) 1.2075.
The origin recorded is \code{matrix}.
\end{ftsevidence}
\end{column}
\end{columns}
\fsisource{\code{workspace/fsi\_setup.py}, \code{resolve\_fsi\_setup};
  \code{docs/fsi-workspace.md}, ``Calibration and provenance''}
\end{frame}

\begin{frame}{A tabulated shear modulus stays tabulated}
\begin{columns}[T]
\begin{column}{0.54\textwidth}
\begin{itemize}\small
  \item A material either takes \(G\) as its source tabulates it, or derives it
        as \(E / (2(1+\nu))\) and says so (\code{shear\_modulus\_basis}).
  \item When the source tabulates \(G\), changing \(E\) or \(\nu\) keeps that
        \(G\). Only a material declaring a derived \(G\) follows the formula.
  \item The two contracts stay distinct even at unity calibration: nothing
        substitutes the isotropic formula for a tabulated value.
  \item Calibration never rewrites the source database; the base and effective
        material both go to the provenance.
\end{itemize}
\end{column}
\begin{column}{0.42\textwidth}
\begin{ftsevidence}[The titanium entry]
\small Tabulated: \(G = 44.0\) GPa.\\
From the same sheet, \(E/(2(1+\nu)) = 42.4\) GPa.\\[1mm]
The entry keeps 44.0, about 4\,\% apart, and a torsion stiffness built from
it keeps it too.
\end{ftsevidence}
\begin{takeaway}
Unity calibration, or a density, \(E\) or \(\nu\) change, does not substitute
the formula for a tabulated \(G\).
\end{takeaway}
\end{column}
\end{columns}
\fsisource{\code{docs/fsi-workspace.md}, ``Calculated solid sections'' and
  ``Calibration and provenance''; \code{fsi/materials.py}}
\end{frame}

\begin{frame}{The pitch factor scales angles, once}
\begin{columns}[T]
\begin{column}{0.54\textwidth}
\begin{itemize}\small
  \item \code{FSI\_GEOMETRIC\_PITCH\_DEG\_FACTOR} scales each signed
        structural angle in degrees once; zero stays zero.
  \item It changes the orientation of the spinning beam's sections used to
        generate the structural nodes and to project the loads; section
        properties, radii and node order stay.
  \item It is not an angular offset. It does not turn the section contours
        again, and it does not re-pitch the aerodynamic mesh.
  \item Keep the aerodynamic geometry consistent with the intended structural
        orientation, and regenerate the node file and ordering map from the
        effective configuration.
\end{itemize}
\end{column}
\begin{column}{0.42\textwidth}
\begin{ftsevidence}[File factor 2.0, row factor 1.5]
\small \code{geometric\_pitch\_deg = [-60, 20]}\\[1mm]
becomes \code{[-90, 30]},\\
not \code{[-180, 60]}.
\end{ftsevidence}
\begin{alertblock}{Alone, it validates nothing}
\small This parameter alone does not establish coupled aerodynamic
validation.
\end{alertblock}
\end{column}
\end{columns}
\fsisource{\code{docs/fsi-workspace.md}, ``Calibration and provenance''}
\end{frame}

\begin{frame}{Provenance: every number says where it came from}
\begin{columns}[T]
\begin{column}{0.50\textwidth}
\textbf{\small \code{fsi-provenance.json}, per point}
\begin{itemize}\scriptsize
  \item the source path and its sha256, and the mode;
  \item the unscaled (base) and effective configurations;
  \item every factor, and whether it came from \code{unity}, \code{file} or
        \code{matrix};
  \item calculated mode: the base and effective material properties and the
        section algorithm's provenance;
  \item \code{section\_model}: \code{"existing solid homogeneous Euler beam"}.
\end{itemize}
\end{column}
\begin{column}{0.46\textwidth}
\textbf{\small Inside a generated blade: \code{BladeProperties.provenance}}
\begin{itemize}\scriptsize
  \item the material entry, its source and the database version;
  \item the geometry file's NAME and sha256, never its path;
  \item the hypotheses, the torsion method and its grid per station, and the
        thin-section ratio per station.
\end{itemize}
\textbf{\small \code{config\_sha256}}
\begin{itemize}\scriptsize
  \item over the sorted, whitespace-free JSON: it identifies the physics, not
        the file layout;
  \item on every row of the convergence log. A typed configuration has no
        provenance and keeps the digest it always had.
\end{itemize}
\end{column}
\end{columns}
\fsisource{\code{workspace/fsi\_setup.py}, \code{ResolvedFsiSetup.provenance};
  \code{fsi/config.py}, \code{config\_sha256}; FSI tutorial, ``config''}
\end{frame}

% Copyright (c) 2026 Geovana Neves. Licensed under CC BY 4.0 (https://creativecommons.org/licenses/by/4.0/). See LICENSE-AND-AUTHORSHIP.md.
%
% fsi/text/05-the-row.tex

\section{From the row to the solver}

\begin{frame}{What a coupled row stages}
\small A row that states \code{FSI: f001} on a workflow that allows it
prepares the driver's complete input interface. Planning derives the nodes
and the maps from the effective configuration and resolves the boundaries and
the frames of each blade or of the wing; the run writes and hashes:
\par\vspace{1.5mm}
{\ftstablesize\renewcommand{\arraystretch}{1.15}
\begin{tabularx}{\linewidth}{@{}P{0.27\linewidth}R@{}}
\toprule
\textbf{File, in the point's directory} & \textbf{What it is} \\
\midrule
\code{config.json}, \code{fsi-provenance.json} & the effective configuration, and where every number came from \\
\code{fsi\_nodes.csv} & the structural nodes, inside the component, in its own frame; one identical file imported once per blade \\
\code{fsi\_node\_map.json} & the node ordering map (\code{node\_map\_file}), from the same generator \\
\code{fsi\_family\_map.json} & the emitted section-load order: which rows of the loads export belong to which blade \\
\code{fsi\_post.txt} & the native loads-export callback: update the sections, compute their loads in Newtons, export them \\
\code{fsi\_callback.py} & the Python callback: one \code{pyfs-fsi} step in the point's directory \\
\bottomrule
\end{tabularx}}
\par\vspace{1.5mm}
{\scriptsize The callback runs under the interpreter that prepared the run
(its \code{pythonw.exe} sibling on Windows): the package, with its
\code{[fsi]} extra, must be installed for that interpreter on the machine
that executes the point.\par}
\fsisource{\code{docs/fsi-workspace.md}, ``Automatic workspace coupling'';
  \code{cases/fsi\_workspace.py}, \code{wire\_workspace\_fsi};
  \code{fsi/beam.py} (PyNite from the \code{[fsi]} extra)}
\end{frame}

% RE-CHECK after the FSI branch (FSI-1, FSI-GUARD, FSI-G) merges: the
% command list and its order are written from the requirements (surface by
% boundary ID, multiquadric kernel, steady and unsteady coupling), not read
% from the merged builder; compare with a planned script of each workflow.
\begin{frame}[fragile]{The native commands, in the order they are emitted}
\relax % a body opening with a brace group would be read as the frame subtitle
{\ftstablesize\renewcommand{\arraystretch}{0.97}
\begin{tabularx}{\linewidth}{@{}P{0.47\linewidth}R@{}}
\toprule
\textbf{Command} & \textbf{What the row gives it} \\
\midrule
\code{DELETE\_AEROELASTIC\_STRUCTURAL\_NODES} & nothing: a clean start \\
\code{AEROELASTIC\_RBF\_TYPE} & \code{MULTI\_QUADRATIC}, for a beam line of nodes \\
\code{ASSIGN\_AEROELASTIC\_SURFACES} & the coupled boundaries, by their identifiers in the mesh \\
\code{ASSIGN\_AEROELASTIC\_COORDINATE\_SYSTEMS} & each blade's own frame, or the wing's \\
\code{IMPORT\_AEROELASTIC\_STRUCTURAL\_NODES} & per blade, in blade order: its frame, \code{DISABLE}, \code{fsi\_nodes.csv} \\
\code{SET\_AEROELASTIC\_WORKING\_DIRECTORY} & \code{.}, the point's directory \\
\code{SET\_AEROELASTIC\_POST\_PROCESSING\_SCRIPT} & \code{fsi\_post.txt} \\
\code{SET\_AEROELASTIC\_STRUCTURAL\_EXECUTION\_COMMAND} & the interpreter and \code{fsi\_callback.py} \\
\code{SET\_AEROELASTIC\_ITERATIONS} & \code{1}, one call per coupling step \\
\code{EXECUTE\_AEROELASTIC\_ANALYSIS} & a \code{steady} or \code{qsteady\_rotor} row: the steady coupled solve \\
\code{SET\_AEROELASTIC\_COUPLING\_IN\_UNSTEADY} & an \code{unsteady} row: \code{ENABLE}, before \code{START\_SOLVER} \\
\bottomrule
\end{tabularx}}
\par\vspace{1mm}
\begin{columns}[T]
\begin{column}{0.44\textwidth}
\scriptsize \code{fsi\_callback.py}, whole:
\begin{lstlisting}[style=ftsbase, language=Python, basicstyle=\ttfamily\tiny]
from pyflightstream.fsi.cli import main
raise SystemExit(main(['step', '--dir', '.']))
\end{lstlisting}
\end{column}
\begin{column}{0.52\textwidth}
\begin{alertblock}{Documented, not verified}
\scriptsize The command database records every command of the family as
documented, none as verified: its 26.124 rows are carried from the 26.123
manual. The coupled runs of this version are measurements, not promotions.
\end{alertblock}
\end{column}
\end{columns}
\fsidecided{\code{cases/fsi\_workspace.py}, \code{wire\_workspace\_fsi};
  \code{commands/aeroelastic\_coupling.yaml}}
\end{frame}

\begin{frame}[fragile]{The pproc sections a coupled row needs}
\begin{columns}[T]
\begin{column}{0.44\textwidth}
\small A one-blade rotor whose family is \code{Blade1}:
\begin{lstlisting}[style=ftsbase, basicstyle=\ttfamily\scriptsize]
[sections]
count = 50
include_symmetry = false

[[sections.distributions]]
families = ["Blade1"]
frame = "LOCAL_AXIS"
planes = ["XY"]
\end{lstlisting}
{\scriptsize No separate \code{[settings.aeroelastic]} table is
needed.\par}
\end{column}
\begin{column}{0.52\textwidth}
\begin{itemize}\small
  \item Exactly one XY distribution per blade, emitted in blade order, each in
        that blade's own frame (span along Z); a wing takes one along its
        span.
  \item Several blades: select the rotor's blades so \code{LOCAL\_AXIS}
        expands them in the declared blade order. The structural
        \code{blade\_count} counts every explicit blade.
  \item Duplicate cuts, aggregated families, other planes or additional
        distributions are refused, never attributed to the first matching
        blade.
\end{itemize}
\begin{takeaway}
The loads export is one flat table in creation order; the family map is what
tells the driver which rows are which blade.
\end{takeaway}
\end{column}
\end{columns}
\fsisource{\code{docs/fsi-workspace.md}, ``Automatic workspace coupling'';
  FSI tutorial, ``loads''}
\end{frame}

% Copyright (c) 2026 Geovana Neves. Licensed under CC BY 4.0 (https://creativecommons.org/licenses/by/4.0/). See LICENSE-AND-AUTHORSHIP.md.
%
% fsi/text/06-refusals.tex
%
% The plan-time refusals are described by what they refuse and why, from the
% requirements of this version; the input-reading messages of the last frame
% are quoted from workspace/fsi_setup.py, shortened only where "..." says so.

\section{What is refused, and why}

\begin{frame}{Refused at plan: the shape of the run}
\small The automatic route is deliberately bounded: each of these refuses the
row when it is planned, before it can run.
\par\vspace{1.5mm}
{\ftstablesize\renewcommand{\arraystretch}{1.18}
\begin{tabularx}{\linewidth}{@{}P{0.30\linewidth}R@{}}
\toprule
\textbf{The row} & \textbf{Why it is refused} \\
\midrule
turns a rotor (\code{unsteady\_rotor}) & the plan says the feature is still in
  debug: a turning blade is deformed in its un-turned position [4] \\
is a \code{qsteady\_rotor} wheel & the coupled quasi-steady case is the sector \\
continues a stopped run & the structural state and the section order of a
  continuation are not proved \\
opens a saved \code{.fsm} & the section order a saved simulation inherits has
  no proof: import the mesh, with its boundary sidecar \\
places nodes in a unit other than metres & other native node units have not
  been measured \\
carries a raw command or a custom flag & the section and frame order can no
  longer be proved \\
puts a structural node outside the component & the interpolation cannot
  place a surface around a node that is not within it \\
\bottomrule
\end{tabularx}}
\fsidecided{\code{cases/fsi\_workspace.py}, \code{validate\_workspace\_fsi};
  \code{cases/workflows.py}}
\end{frame}

\begin{frame}{Refused at plan: blades, frames and the speed}
\begin{itemize}\small
  \item \textbf{The count.} The structural \code{blade\_count} must equal the
        rotor's blades, every one explicit.
  \item \textbf{The speed.} \(\Omega = |\mathrm{rpm}|\,\pi/30\) must equal the
        row's rotor speed; on a \code{qsteady\_rotor} row that is the speed of
        its free-stream rotation. A structure that does not turn has
        \(\Omega = 0\) and carries its own weight instead.
  \item \textbf{The clock.} On an \code{unsteady} row the structural time
        increment must equal the solver's time step.
  \item \textbf{The mapping.} Blade, frame and boundary are one-to-one; there
        is exactly one section distribution per blade, in blade import order,
        each in its own frame with XY cuts, span along Z.
\end{itemize}
{\scriptsize Geometry and structural pitch must describe the same blade: the
adapter infers no basis and re-pitches no mesh.\par}
\fsidecided{\code{cases/fsi\_workspace.py}, \code{wire\_workspace\_fsi},
  \code{validate\_workspace\_fsi}}
\end{frame}

\begin{frame}{Refused while reading the input}
\relax % a body opening with a brace group would be read as the frame subtitle
{\ftstablesize\renewcommand{\arraystretch}{1.18}
\begin{tabularx}{\linewidth}{@{}P{0.28\linewidth}R@{}}
\toprule
\textbf{The input} & \textbf{The message says} \\
\midrule
a factor key the package does not know & \says{unknown FSI matrix keys [...]; use FSI\_GEOMETRIC\_PITCH\_DEG\_FACTOR, ...} \\
factors on a row with no \code{FSI} & \says{matrix calibration states no FSI input; add FSI: f001 in VAR\_NAMES\_VALUES.} \\
a code that is not \code{f} + name & \says{FSI='...'; name an f-prefixed code without an extension, such as f001.} \\
a zero or negative factor & \says{... use a finite positive factor.} \\
a stiffness scale other than 1 in \code{[config]} & \says{... config.stiffness\_scale\_factor must be 1.0 in workspace inputs; calibrate bending\_stiffness\_n\_m2 and torsion\_stiffness\_n\_m2 explicitly ...} \\
material factors in supplied mode & \says{supplied mode cannot apply material factors [...] to already supplied distributions; calibrate the distribution itself.} \\
a blade table in calculated mode & \says{calculated mode derives config.blade; remove the supplied blade table.} \\
a material the database lacks & \says{no material '...' in the materials database (version 1); it holds: al-7075-t6, ti-6al-4v-grade5-annealed. ...} \\
\bottomrule
\end{tabularx}}
\fsisource{\code{workspace/fsi\_setup.py}, \code{resolve\_row\_fsi},
  \code{resolve\_fsi\_setup}; \code{fsi/materials.py}, \code{material}}
\end{frame}

% Copyright (c) 2026 Geovana Neves. Licensed under CC BY 4.0 (https://creativecommons.org/licenses/by/4.0/). See LICENSE-AND-AUTHORSHIP.md.
%
% fsi/text/07-examples.tex
%
% The numbers are what the three examples print when run against the tag's
% sources (python examples/<name>.py, PYTHONPATH at the tag's src/, the [fsi]
% extra installed), on 2026-09-28. Every blade is synthetic by the package's
% policy. A release that changes an example changes these numbers: run it
% again and copy them.

\section{Worked examples}

\begin{frame}[fragile]{Example 1: calculated and supplied inputs agree}
\begin{columns}[T]
\begin{column}{0.54\textwidth}
\code{examples/workspace\_fsi\_calibration.py}, offline, no solver:
\begin{enumerate}\scriptsize
  \item Writes \code{inputs/fsi/f\_calculated.toml} from the package template
        (the titanium strip), with the file factor 1.05 on EI.
  \item Resolves it and reads the \emph{unscaled} base distributions.
  \item Writes \code{f\_supplied.toml} with those base distributions and the
        same file factor 1.05. Reusing the effective values here would apply
        the calibration a second time.
  \item Resolves both with \code{FSI\_BENDING\_STIFFNESS\_N\_M2\_FACTOR: 1.15}.
  \item Asserts every distribution equal across the two modes, effective EI
        equal to base EI times 1.15, and the origin \code{matrix}.
\end{enumerate}
\begin{lstlisting}[style=ftsbase, basicstyle=\ttfamily\tiny]
python workspace_fsi_calibration.py /absolute/path/to/new-workspace
\end{lstlisting}
\end{column}
\begin{column}{0.42\textwidth}
{\ftstablesize\renewcommand{\arraystretch}{1.2}
\begin{tabular}{@{}lr@{}}
\toprule
\textbf{Printed} & \textbf{Value} \\
\midrule
base EI, each station & 24.277 N\,m\(^2\) \\
effective EI, each station & 27.919 N\,m\(^2\) \\
file factor, matrix factor & 1.05, 1.15 \\
\code{factor\_origin} & \code{matrix} \\
\code{mode\_equivalence} & verified \\
\code{solver\_executed} & \code{false} \\
\bottomrule
\end{tabular}}
\begin{ftsevidence}[Check it by hand]
\scriptsize \(I = b h^3/12 = 0.04 \times 0.004^3 / 12 = 2.133\times10^{-10}\) m\(^4\);
\(E I = 113.8\) GPa \(\times\, I = 24.28\) N\,m\(^2\);
\(\times\, 1.15 = 27.92\), not \(\times\, 1.05 \times 1.15\).
\end{ftsevidence}
{\scriptsize Choose a new destination: the example refuses to overwrite its
two inputs.\par}
\end{column}
\end{columns}
\fsisource{\code{examples/workspace\_fsi\_calibration.py}, run at the tag;
  \code{workspace/fsi\_setup.py}, \code{FSI\_TEMPLATE}}
\end{frame}

\begin{frame}{Example 2: one section, and its torsion constant}
\begin{columns}[T]
\begin{column}{0.48\textwidth}
\code{examples/fsi\_solid\_blade\_properties.py} builds a synthetic blade:
seven stations from 0.15 to 0.75 m, chord 80 to 40 mm, pitch 35 to 11 deg,
NACA 2415 at the root thinning to NACA 2409 at the tip, written to a section
file and read back, the file's sha256 into the provenance.
\par\vspace{2mm}
{\tiny\renewcommand{\arraystretch}{1.15}
\begin{tabular}{@{}lr@{}}
\toprule
\textbf{Root section, NACA 2415} & \\
\midrule
area & 654.144 mm\(^2\) \\
centroid from the pitch axis & \(-13.422\) mm chordwise \\
\(I\) flap, \(I\) chord & 5512.2, 227784.7 mm\(^4\) \\
principal angle & \(-0.100\) deg \\
\bottomrule
\end{tabular}}
\end{column}
\begin{column}{0.48\textwidth}
{\ftstablesize\renewcommand{\arraystretch}{1.15}
\begin{tabular}{@{}lr@{}}
\toprule
\textbf{Grid, cells across} & \(J\) [mm\(^4\)] \\
\midrule
16 & 20926.02 \\
32 & 20956.81 \\
64, the default & 20962.52 \\
thin-section estimate & 21721.14 \\
\bottomrule
\end{tabular}}
\par\vspace{1.5mm}
\begin{itemize}\scriptsize
  \item Doubling the cells divides the error by about four: at 64 cells about
        a third of the change from 32 to 64 is left.
  \item The thin-section formula is 3.6\,\% high here: a cross-check, never
        the value.
\end{itemize}
\end{column}
\end{columns}
\fsisource{\code{examples/fsi\_solid\_blade\_properties.py}, run at the tag}
\end{frame}

\begin{frame}{Example 2: the whole blade, in titanium}
\relax % a body opening with a brace group would be read as the frame subtitle
{\ftstablesize\renewcommand{\arraystretch}{1.1}
\begin{tabular*}{\linewidth}{@{\extracolsep{\fill}}rrrrrrrr@{}}
\toprule
\(r\) [m] & \(\mu\) [kg/m] & EI [N\,m\(^2\)] & GJ [N\,m\(^2\)] & \(I_1\) [kg\,m] & \(I_2\) [kg\,m] & \(e_c\) [mm] & thin\(/J\) \\
\midrule
0.15 & 2.8979 & 627.29 & 922.35 & 1.009e-03 & 2.442e-05 & \(-13.42\) & 1.036 \\
0.25 & 2.2727 & 361.01 & 531.76 & 6.650e-04 & 1.405e-05 & \(-12.30\) & 1.032 \\
0.35 & 1.7441 & 198.04 & 291.97 & 4.217e-04 & 7.708e-06 & \(-11.19\) & 1.028 \\
0.45 & 1.3040 & 102.60 & 151.24 & 2.554e-04 & 3.993e-06 & \(-10.07\) & 1.024 \\
0.55 & 0.9445 & 49.58 & 72.98 & 1.462e-04 & 1.930e-06 & \(-8.95\) & 1.020 \\
0.65 & 0.6574 & 21.98 & 32.24 & 7.789e-05 & 8.553e-07 & \(-7.83\) & 1.017 \\
0.75 & 0.4347 & 8.73 & 12.72 & 3.784e-05 & 3.395e-07 & \(-6.71\) & 1.014 \\
\bottomrule
\end{tabular*}}
\par\vspace{2mm}
\begin{columns}[T]
\begin{column}{0.56\textwidth}
\begin{itemize}\scriptsize
  \item Ti-6Al-4V grade 5, annealed: \(\rho\) 4430 kg/m\(^3\), \(E\) 113.8 GPa,
        \(G\) 44.0 GPa, \(\nu\) 0.342, database version 1.
  \item The configuration round-trips through \code{config.json} unchanged,
        with its provenance and its \code{config\_sha256}.
  \item \(e_c\): the centroid ahead of or behind the pitch axis, which the
        elastic axis is assumed to follow.
\end{itemize}
\end{column}
\begin{column}{0.40\textwidth}
\begin{ftsevidence}[At rest, from the beam]
\small first flap: 30.6 Hz\\
first torsion: 601.9 Hz
\end{ftsevidence}
\end{column}
\end{columns}
\fsisource{\code{examples/fsi\_solid\_blade\_properties.py}, sections 3 and 4,
  run at the tag}
\end{frame}

\begin{frame}{Example 3: the Campbell diagram of a rotating blade}
\begin{columns}[T]
\begin{column}{0.58\textwidth}
\begin{center}
\begin{adjustbox}{max width=\linewidth, max totalheight=0.70\textheight}
\begin{tikzpicture}[x=0.085cm, y=0.0118cm]
  % axes
  \draw[->, SoftGray] (0,0) -- (98,0);
  \node[font=\tiny, anchor=north] at (48,-30) {rotor speed \(\Omega\) [rad/s]};
  \draw[->, SoftGray] (0,0) -- (0,430);
  \node[font=\tiny, anchor=south, rotate=90] at (-14,215) {natural frequency \(\omega_n\) [rad/s]};
  \foreach \x in {0,15,30,45,60,75,90}
    \draw[SoftGray] (\x,0) -- (\x,-6) node[below, font=\tiny, text=black] {\x};
  \foreach \y in {0,100,200,300,400}
    \draw[SoftGray] (0,\y) -- (-1.5,\y) node[left, font=\tiny, text=black] {\y};
  % legend, in the empty upper left
  \draw[thick, CourseBlue] (5,405) -- (13,405);
  \fill[CourseBlue] (9,405) circle (1.2pt);
  \node[font=\tiny, anchor=west] at (14,405) {first flap};
  \draw[thick, CourseRed] (5,378) -- (13,378);
  \fill[CourseRed] (8.2,375.5) rectangle (9.8,380.5);
  \node[font=\tiny, anchor=west] at (14,378) {first torsion};
  \draw[dotted, SoftGray, thick] (5,351) -- (13,351);
  \node[font=\tiny, anchor=west] at (14,351) {\(n\)P excitation, \(n\,\Omega\)};
  % excitation rays nP
  \foreach \n in {1,2,3,4} {
    \draw[dotted, SoftGray, thick] (0,0) -- (90,{\n*90});
    \node[font=\tiny, text=SoftGray, right] at (90,{\n*90}) {\n P};
  }
  % first flap
  \draw[thick, CourseBlue] plot[mark=*, mark size=1.2pt] coordinates
    {(0,27.17) (15,31.70) (30,42.37) (45,55.55) (60,69.69) (75,84.23) (90,98.95)};
  % first torsion
  \draw[thick, CourseRed] plot[mark=square*, mark size=1.2pt] coordinates
    {(0,310.90) (15,311.25) (30,312.29) (45,314.02) (60,316.42) (75,319.48) (90,323.18)};
\end{tikzpicture}
\end{adjustbox}
\end{center}
\end{column}
\begin{column}{0.38\textwidth}
{\scriptsize A uniform synthetic blade, 1 m, 15 stations, swept from 0 to 90
rad/s. At every speed: static solve under its own centrifugal loads, then the
modes with the tension's geometric stiffness.\par}
\vspace{1.5mm}
{\ftstablesize\renewcommand{\arraystretch}{1.15}
\begin{tabular}{@{}lrr@{}}
\toprule
\textbf{Southwell fit} & \textbf{flap} & \textbf{torsion} \\
\midrule
\(\omega_0\) [rad/s] & 27.96 & 310.90 \\
\(S\) & 1.118 & 0.961 \\
\(r^2\) & 0.99986 & 1.00000 \\
\bottomrule
\end{tabular}}
\par\vspace{1.5mm}
{\scriptsize \(\omega_n^2 = \omega_0^2 + S\,\Omega^2\). Flap \(S\) sits in
the expected 1.1 to 1.3; torsion near \((I_1 - I_2)/(I_1 + I_2) = 0.961\).
A crossing of a track and a ray flags a potential resonance [10, 11].\par}
\end{column}
\end{columns}
\fsisource{\code{examples/fsi\_campbell\_diagram.py}, run at the tag;
  FSI tutorial, ``centrifugal''}
\end{frame}

% Copyright (c) 2026 Geovana Neves. Licensed under CC BY 4.0 (https://creativecommons.org/licenses/by/4.0/). See LICENSE-AND-AUTHORSHIP.md.
%
% fsi/text/08-outputs-and-limits.tex
%
% The two limits frames state what the sources say is NOT established. They
% are the frames never to soften: a later release moves a sentence here only
% when its own documentation moves it.

\section{Outputs, checklist and limits}

\begin{frame}{Outputs, all in the point's working directory}
\relax % a body opening with a brace group would be read as the frame subtitle
{\ftstablesize\renewcommand{\arraystretch}{1.1}
\begin{tabularx}{\linewidth}{@{}P{0.36\linewidth}R@{}}
\toprule
\textbf{File, written at every call} & \textbf{What it holds} \\
\midrule
\code{FS\_SurfaceSection\_Loads.txt} & by the solver: the section loads, computed in Newtons, line densities along the span; the reader asserts units and completeness \\
\code{FSIDisp.txt} & one translation per node, in metres, blade frame, import order \\
\code{state.json} & the structural state, written atomically, so a killed call resumes \\
\code{fsi\_convergence\_log.csv} & the validity boundary in its header, \code{config\_sha256} on every row \\
\code{fsi\_archive/call\_NNNN/} & that call's loads and displacements, for offline replay \\
\code{pyfs\_fsi\_calls.log}, \code{\_error.log} & a line per call; on failure the traceback, and a non-zero exit that stops the coupling rather than let it continue rigid \\
\bottomrule
\end{tabularx}}
\par\vspace{1mm}
\begin{itemize}\scriptsize
  \item \code{pyfs-fsi step --dir /absolute/path} runs one driver call when
        current loads and maps already exist; stale solver iterations are refused.
  \item Changing a calibration requires a new resolved point and regenerated
        nodes.
\end{itemize}
\fsisource{\code{docs/fsi-workspace.md}, ``Automatic workspace coupling'';
  \code{fsi/cli.py}; \code{fsi/driver.py}; FSI tutorial, ``loads'', ``driver
  and state'', ``cli''}
\end{frame}

\begin{frame}{Before planning an FSI row}
\begin{enumerate}\small
  \item \textbf{The input}: \code{inputs/fsi/f<id>.toml} in ONE mode, SI units;
        calculated mode names a database material or a fully sourced one, and
        closed simple contours in metres about the pitch axis.
  \item \textbf{The row}: \code{WORKFLOW} \code{steady}, \code{unsteady}, or
        \code{qsteady\_rotor} with periodic symmetry (the sector);
        \code{FSI: f<id>} and any \code{FSI\_*\_FACTOR} in
        \code{VAR\_NAMES\_VALUES}; no \code{RAW}, no custom flags.
  \item \textbf{The geometry}: a fresh OBJ or STL import with its boundary
        sidecar, in \code{METER}; never a saved \code{.fsm}.
  \item \textbf{The rotor}: one isolated, axisymmetric rotor; shaft X, blade
        datum Z; \code{blade\_count} equal to its blades;
        \code{omega\_rad\_per\_s} equal to the row's free-stream rotation. A
        wing: \code{[config.wing]}, \(\Omega = 0\), \code{blade\_count = 1}.
  \item \textbf{The pproc}: one XY distribution per blade, \code{LOCAL\_AXIS},
        in blade order, \code{include\_symmetry = false}.
  \item \textbf{Then} \code{pyfs-matrix plan}, where the refusals of this deck
        fire; after the run, \code{fsi-provenance.json} in the point's directory
        gives each factor's value and origin.
\end{enumerate}
\fsisource{\code{docs/fsi-workspace.md}; \code{cases/fsi\_workspace.py};
  \code{workspace/fsi\_setup.py}}
\end{frame}

\begin{frame}{Limits: what the coupling does not establish}
\begin{itemize}\small
  \item \textbf{Wiring is not validation}: the staging and callback checks do
        not establish native two-way rotor deformation, convergence or
        aerodynamic validity.
  \item \textbf{The native evidence so far} (RPT-093 [2]): on 26.124 a blade
        held still in a turning free stream, steady, mapped every surface
        vertex and exported the imposed flap to 0.1 mm, its rigid axial force
        within 0.35\,\% of the same blade turning; a fixed wing mapped and
        exported the bend, steady and unsteady. A turning blade is morphed at
        its imported azimuth, which is why \code{unsteady\_rotor} with FSI is
        refused. These are imposed-displacement probes, not coupled accuracy.
  \item \textbf{Open}: saved-simulation and continuation use, and native
        coupled validation; staging or command acceptance implies neither.
  \item \textbf{The zero-speed example inputs} are for offline checks, not
        runnable coupled rotor rows.
\end{itemize}
\begin{ftslimit}[The package's own capability table]
\scriptsize FSI beam and modal analysis: experimental. The coupled driver:
experimental, replayed offline on archived fixtures, never run against a live
solver in CI. Rotary two-way coupling: not validated.
\end{ftslimit}
\vspace{-1.5mm}
\fsisource{\code{docs/fsi-workspace.md}, ``Automatic workspace coupling'';
  \code{README.md}, the capability status table;
  \code{docs/srs/architecture-srs.md}}
\end{frame}

\begin{frame}{Limits: the structural model}
\begin{itemize}\small
  \item \textbf{An Euler beam}, one per blade, in flap and torsion
        \((w, \theta)\): only those two carry mass. This input adds no shell,
        spar, hollow-section, shear-centre or second structural model; hollow
        and spar sections are a future option, not built.
  \item \textbf{Solid, homogeneous sections} from ONE material, when
        calculated.
  \item \textbf{The elastic axis at the centroid}; the shear centre is not
        computed. A generated blade has a zero CG offset, so the bend-twist
        coupling that offset drives is absent.
  \item \textbf{No trapeze effect}: the elastic twist is overpredicted, which
        is conservative.
  \item \textbf{Quasi-static}: the once-per-revolution content is trustworthy
        only well below resonance, \(n\,\Omega / \omega_n \lesssim 0.3\).
  \item \textbf{Sources}: the formulas are transcribed from the coupling plan,
        not re-checked against the primary papers; the materials are typical
        values, not design allowables.
\end{itemize}
\fsisource{\code{docs/fsi-workspace.md}; FSI tutorial, ``The loop in one
  paragraph'', ``sections and materials'', ``beam'', ``centrifugal'';
  \code{fsi/materials.py}}
\end{frame}


\begin{frame}{The physics behind the coupling, and where to read it}
\begin{itemize}\small
  \item \textbf{The beam}: Euler bending and Saint-Venant torsion of a
        slender member, and the section constants of a solid contour [6, 7].
  \item \textbf{The rotating blade}: centrifugal tension and stiffening, the
        propeller moment toward flat pitch, the in-plane softening, and the
        natural frequencies that rise with the speed [8, 9, 10].
  \item \textbf{The Campbell diagram}: the blade's frequencies against the
        rotor speed, beside the lines of the rotor's harmonics [11].
  \item \textbf{Quasi-static coupling}: the unsteady aerodynamics of an
        oscillating section, and the reduced frequency below which a steady
        answer holds [12, 13].
  \item \textbf{The interpolation}: radial basis functions from a set of
        structural nodes to an aerodynamic surface, and the multiquadric
        kernel [14, 15, 16].
\end{itemize}
\end{frame}

\begin{frame}{References}
\begin{reflist}
  \bibitem{g06r01} \docref{fsi-workspace}{FSI in a workspace}; and the FSI tutorial,
        \code{src/pyflightstream/fsi/README.md} of the repository.
  \bibitem{g06r02} \pkgauthor, ``The Aeroelastic Coupling Toolbox on 26.124: nine
        measured facts'', report RPT-093, 2026, \code{reports/} of
        \url{\pkgrepo}.
  \bibitem{g06r03} \pkgauthor, ``The quasi-steady rotor against the unsteady rotor,
        measured'', report RPT-089, 2026, \code{reports/} of \url{\pkgrepo}.
  \bibitem{g06r04} \recordsref
  \bibitem{g06r05} D. C. Brinck, \emph{PyNite}: a 3D structural finite element library
        for Python. \url{https://github.com/JWock82/Pynite}
  \bibitem{g06r06} S. P. Timoshenko, J. N. Goodier, \emph{Theory of Elasticity}, 3rd ed.,
        McGraw-Hill, New York, 1970.
  \bibitem{g06r07} W. D. Pilkey, \emph{Analysis and Design of Elastic Beams:
        Computational Methods}, Wiley, New York, 2002.
  \bibitem{g06r08} D. H. Hodges, E. H. Dowell, ``Nonlinear equations of motion for the
        elastic bending and torsion of twisted nonuniform rotor blades'',
        NASA TN D-7818, 1974.
\end{reflist}
\end{frame}

\begin{frame}{References (continued)}
\begin{reflist}
  \setcounter{enumi}{8}
  \bibitem{g06r09} R. L. Bielawa, \emph{Rotary Wing Structural Dynamics and
        Aeroelasticity}, AIAA Education Series, AIAA, Washington, 1992.
  \bibitem{g06r10} R. V. Southwell, B. S. Gough, ``On the free transverse vibration of
        airscrew blades'', Aeronautical Research Committee, Reports and
        Memoranda 766, 1921.
  \bibitem{g06r11} W. Campbell, ``The protection of steam-turbine disk wheels from axial
        vibration'', \emph{Transactions of the ASME}, 46, pp. 31-160, 1924.
  \bibitem{g06r12} R. L. Bisplinghoff, H. Ashley, R. L. Halfman, \emph{Aeroelasticity},
        Addison-Wesley, Reading, 1955.
  \bibitem{g06r13} T. Theodorsen, ``General theory of aerodynamic instability and the
        mechanism of flutter'', NACA Report 496, 1935.
  \bibitem{g06r14} R. L. Hardy, ``Multiquadric equations of topography and other
        irregular surfaces'', \emph{Journal of Geophysical Research}, 76(8),
        pp. 1905-1915, 1971.
  \bibitem{g06r15} A. Beckert, H. Wendland, ``Multivariate interpolation for
        fluid-structure-interaction problems using radial basis functions'',
        \emph{Aerospace Science and Technology}, 5(2), pp. 125-134, 2001.
  \bibitem{g06r16} A. de Boer, M. S. van der Schoot, H. Bijl, ``Mesh deformation based
        on radial basis function interpolation'', \emph{Computers and
        Structures}, 85(11-14), pp. 784-795, 2007.
\end{reflist}
\end{frame}

\end{document}
